\begin{frame}
    \frametitle{Ejercicio}

    Estudia la dependencia o independencia lineal de los siguientes vectores
    de \(\mathbb{R}^3\):
    \[
        v_1=(1,2,1),\qquad
        v_2=(2,1,0),\qquad
        v_3=(3,3,1).
    \]

    Si son linealmente dependientes, encuentra una relación de dependencia
    lineal.
\end{frame}


\begin{frame}
    \frametitle{Ejercicio}

    Determina para qué valores de \(a\in\mathbb{R}\) son linealmente
    independientes los vectores
    \[
        v_1=(1,0,1),\qquad
        v_2=(0,1,1),\qquad
        v_3=(1,1,a).
    \]

    Para los restantes valores de \(a\), encuentra una relación de
    dependencia lineal.
\end{frame}


\begin{frame}
    \frametitle{Ejercicio}

    Estudia la dependencia o independencia lineal en \(\mathbb{R}_2[x]\)
    de los polinomios
    \[
        p_1(x)=1+x,
        \qquad
        p_2(x)=1-x,
        \qquad
        p_3(x)=1+x^2.
    \]

    \medskip

    Para ello, estudia cuándo
    \[
        \lambda_1p_1(x)+
        \lambda_2p_2(x)+
        \lambda_3p_3(x)=0.
    \]
\end{frame}


\begin{frame}
    \frametitle{Ejercicio}

    En \(M_{2\times2}(\mathbb{R})\), consideremos
    \[
        A_1=
        \begin{pmatrix}
            1&0\\
            0&1
        \end{pmatrix},
        \qquad
        A_2=
        \begin{pmatrix}
            0&1\\
            1&0
        \end{pmatrix},
        \qquad
        A_3=
        \begin{pmatrix}
            1&1\\
            1&1
        \end{pmatrix}.
    \]

    Estudia si \(A_1,A_2,A_3\) son linealmente dependientes.

    Si lo son, encuentra una relación de dependencia lineal.
\end{frame}
